% Propietario conservado: /Users/ruben/Documents/New project/output/EXTREMA_MEDIA_RAZON_RECIPROCIDAD_ES_EN_20260918/ARTICULO/ES/source/sections/registro_k.tex
\section{Extraction dodécaphasique de la constante de couplage arithmético-géométrique}
\label{x_reciprocidad:sec:k-registro}

\begin{definition}[Constante holographique de couplage arithmético-géométrique]
\label{x_reciprocidad:lib01:definicion}
On dénomme ainsi le registre fini ordonné produit par l’extraction
orientée de l’état APP--TRIT--TPK, avec son origine de phase, sa normalisation et
ses lecteurs déclarés, dont la représentation positionnelle intervient dans la
clôture arithmétique et dont les réalisations incidentielle et géométrique agissent
sur le même registre. Le symbole \(K\) désigne sa représentation
entière canonique ; \(\kappa_{\rm ag}\) désigne sa lecture rationnelle
périodique dans la carte fixée.
\end{definition}

La génération des coordonnées régionales n'épuise pas l'information de l'état APP--TRIT--TPK. Une lecture archimédienne détermine une valeur au moyen de cylindres compatibles ; une autre lecture conserve l'opposition des feuillets, l'orientation de la route et le registre des incidences traversées. Le registre dodécaphasique appartient à cette seconde lecture avant d'intervenir dans la clôture d'alpha. Il n'est donc pas défini en soustrayant une constante physique à une combinaison de valeurs préalablement connues. Il est reconstruit à partir d'une coordonnée signée du même état arithmétique avec mémoire de trajectoire qui produit les coordonnées régionales.

La distinction est nécessaire même lorsque deux objets possèdent douze
coordonnées. Un mot de douze blocs du catalogue, une liste de douze
événements orientés, un vecteur entier signé et un registre dodécaphasique en base mille sont des
objets distincts. Leur comparaison requiert les applications qui les
relient. L'égalité de longueur ne fournit pas ces applications, et la
perte des signes n'est pas une simple simplification de notation.

Ce chapitre maintient la délimitation causale fixée dans les chapitres précédents :
les neuf positions APP des axes horizontal et vertical produisent les
feuillets arithmétiques ; TRIT type le régime et l'orientation ; TPK sélectionne,
transporte, met à jour et prolonge, en conservant le résidu, le quotient, la retenue,
le feuillet, la route, la frontière et la mémoire. À partir de la structure discrète conjointe du
continu s'obtiennent les coordonnées utilisées ici. La chaîne
\(w_6\to w_{12}\to w_{18}\to w_{24}\to w_{30}\to R_{36}\to G_9\)
n'est pas remplacée par un calendrier fini. L'ordre
\(U_t\to\Gamma_9\to K_{\mathrm{ph}}\) n'est pas non plus inversé : la dernière application est une
lecture à mémoire réduite, et non le générateur du registre signé.


% BEGIN OWNER /Users/ruben/Documents/New project/output/EXTREMA_MEDIA_RAZON_RECIPROCIDAD_ES_EN_20260918/ARTICULO/ES/source/sections/revision_k.tex
% Adición propuesta al inicio de registro_k.tex, después de su introducción.
% Conserva todas las definiciones, fórmulas y demostraciones actuales.
% Procedencia: registro_k.tex, §§k-hadamard/k-transversal/k-kappas;
% k_reversibilidad.tex; incidencia_articulacion.tex, inc-ampl-cuaterna.
\paragraph{Coordonnées réversibles et détermination incidencielle.}
\label{x_reciprocidad:ampl-alfa-k-intro-reversible}
Le résultat qui organise ce chapitre est l'équivalence entre trois
systèmes de coordonnées d'un registre préalablement produit par
APP--TRIT--TPK. La transformation de Hadamard relie le registre
signé au registre entier ; l'extracteur de différences relie
ce dernier à ses variations transversales et à sa composante uniforme :
\[
 U\underset{\mathcal H_{12}}{\overset{\mathcal H_{12}/4}{\longleftrightarrow}}K
 \overset{(D_3,D_4,Q)}{\longleftrightarrow}
 (D_3K,D_4K,Q(K)).
\]
La première équivalence s'établit sur le sous-réseau entier
$\mathcal L_H$ ; la seconde, sur l'image de l'extracteur. Une fois fixés la
base $1000$, les $12$ positions et l'origine de phase, l'évaluation
$K\mapsto\kappa_{\mathrm{per}}$ ajoute un quatrième système de coordonnées
réversible : la multiplication par $1000^{12}-1$ suivie des divisions
euclidiennes restitue le registre complet. Les démonstrations de
\S\ref{x_reciprocidad:subsec:k-hadamard}, \S\ref{x_reciprocidad:subsec:k-transversal} et
\S\ref{x_reciprocidad:subsec:k-kappas} établissent ces équivalences dans leurs domaines
respectifs.

Deux compositions agissent ensuite sur ce registre. La composition
positionnelle combine ses composantes avec les coordonnées de clôture,
de propagation et d'autoéchelle, et sa normalisation détermine $\alpha_A$. La
composition incidencielle applique le seuil cubique et obtient
\[
 K\longmapsto B_K=\{j:K_j\geq729\}=\{4,6,9,10\}.
\]
Le registre intégral permet de reconstruire les deux compositions. Le support
$B_K$ conserve, spécifiquement, l'appartenance à la classe sélectionnée
par le seuil. L'identité $B_K=H_e\cap P_\pi$ de
\eqref{x_reciprocidad:inc-ampl-cuaterna} relie cette extraction entière à l'incidence
des codes régionaux. Le support négatif de leur bilan avec $K$
complète le drapeau $B_K\subset H^-_\alpha$ développé dans le
prolongement exceptionnel. Ainsi, la fonction de $K$ est définie par
une transformation inversible et par des applications ultérieures
précises ; l'information supplémentaire de profondeur et de transport demeure
dans l'état dont provient le registre.

% END OWNER


\subsection{Extension cyclique et graduation de mémoire}
\label{x_reciprocidad:subsec:k-calendario}

Le calendrier observable contient cent huit positions, organisées en douze fenêtres nonadiques. Avec des indices positifs, soient
\[
 E_m=\{9(m-1)+1,\ldots,9m\},\qquad 1\leq m\leq12.
\]
Le support ordonné \(E_1\sqcup\cdots\sqcup E_{12}\) conserve la position et l'
appartenance à une fenêtre. Le groupe \(C_{108}=\mathbb Z/108\mathbb Z\)
conserve, en revanche, la loi de translation cyclique. Pour le représentant
\(0\leq\widetilde\theta<108\), les coordonnées
\[
 \beta(\theta)=\left[\left\lfloor\frac{\widetilde\theta}{9}\right\rfloor\right]_{12},
 \qquad r(\theta)=1+(\widetilde\theta\bmod9)
\]
séparent le secteur dodécaphasique et la position nonadique intérieure. Aucune d'elles ne conserve à elle seule le nombre total de tours.

\begin{proposition}[La retenue du calendrier]
\label{x_reciprocidad:prop:k-extension}
Les applications
\[
 \iota:C_{12}\longrightarrow C_{108},\quad[b]\longmapsto[9b],
 \qquad p:C_{108}\longrightarrow C_9,\quad[\theta]\longmapsto[\theta]_9
\]
forment une suite exacte non scindée
\[
 0\longrightarrow C_{12}\xrightarrow{\iota}C_{108}
 \xrightarrow{p}C_9\longrightarrow0.
\]
Pour la section ensembliste qui choisit les représentants \(0,\ldots,8\), le défaut d'additivité est la retenue
\[
 c(r,r')=\left[\left\lfloor\frac{\widetilde r+\widetilde r'}9\right\rfloor\right]_{12}.
\]
\end{proposition}

\begin{proof}
Le noyau de \(p\) est formé des multiples de neuf et coïncide avec
l'image de \(\iota\). La réduction est surjective et la multiplication
par neuf est injective sur le domaine indiqué. Si la suite était
scindée, le groupe central serait isomorphe à \(C_{12}\times C_9\), dont l'
exposant est \(\operatorname{mcm}(12,9)=36\) ; l'exposant de
\(C_{108}\) est 108. La contradiction prouve la non-scission. Enfin,
la division euclidienne de \(\widetilde r+\widetilde r'\) par neuf donne la
retenue écrite et vérifie l'identité de la section.
\end{proof}

Le résultat explique une différence qui réapparaîtra dans alpha : revenir à la
phase n’équivaut pas à revenir à l’état. Si \(q_{\mathrm{mem}}\) compte les
tours nonadiques, le calendrier conserve seulement
\(\beta=[q_{\mathrm{mem}}]_{12}\). Après cent huit étapes,
\(\beta\) revient, mais \(q_{\mathrm{mem}}\) augmente de douze. Le
registre conserve en outre le cylindre, le registre des incidences et la mémoire
verticale. La fermeture d’une carte finie ne transforme pas cette histoire en un
cycle sans mémoire.

\subsection{Le registre orienté et la carte intégrale \texorpdfstring{\(1+5+6\)}{1+5+6}}
\label{x_reciprocidad:subsec:k-registro-integral}

Pour une histoire produite par TPK, soit \(\tau_m\) le sous-chemin sur \(E_m\), \(\sigma_m\) son orientation, \(\kappa_m\) la retenue de frontière et \(\Lambda_m\) l'article local d'incidences. Le registre est
\[
 \mathcal R_{12}(x)=
 \bigl((E_m,\tau_m,\sigma_m,\kappa_m,\Lambda_m)\bigr)_{m=1}^{12}.
\]
Son domaine est l'espace des histoires arithmétiques avec mémoire de trajectoire, et non l'ensemble des
étiquettes de phase. La projection sur le calendrier oublie précisément certaines
des données dont le lecteur signé a besoin.

Une carte initiale du registre peut s'écrire au moyen d'événements orientés. L'extraction orientée fixe l'ensemble des codes d'événement \(\{2,4,6,8\}\) et l'orientation sectorielle
\[
 \Sigma_{12}=(+,+,+,-,-,-,+,+,+,-,-,-).
\]
Si \(d_t\) est le code de direction au pas \(t\), on obtient
\[
 e_i=\Sigma_{12}(i+1)
 \#\{t\in E_{i+1}:d_t\in\{2,4,6,8\}\},
 \qquad0\leq i<12.
\]
La formule maintient distincts le recensement non négatif des événements et
l'orientation. Elle n'identifie pas une direction cardinale à un signe TRIT. Le
vecteur \(e\), dont les composantes sont des recensements locaux signés, n'est pas non plus
encore le vecteur \(U\) qui apparaîtra plus bas.

Les positions opposées déterminent, pour \(0\leq i<6\),
\[
 p_i=e_i+e_{i+6},\qquad a_i=e_i-e_{i+6}.
\]
Avec
\[
 s_C=\sum_{i=0}^{5}p_i,\qquad M_i=p_i-p_5\quad(0\leq i<5),
\]
on définit la carte entière
\[
 \mathcal L_{12}(e)
 =(s_C,M_0,\ldots,M_4,a_0,\ldots,a_5).
\]
La décomposition \(1+5+6\) n'est pas une répartition arbitraire des coordonnées :
une coordonnée conserve la charge, cinq comparent les paires par rapport à une
paire de référence et six conservent l'opposition des demi-couronnes.

\begin{lemma}[Inversion de la carte entière]
\label{x_reciprocidad:lem:k-carta-integral}
Un tuple entier \((s_C,M_0,\ldots,M_4,a_0,\ldots,a_5)\) appartient à l'image de \(\mathcal L_{12}\) si et seulement si
\[
 s_C-\sum_{i=0}^{4}M_i\equiv0\pmod6,
 \qquad p_i\equiv a_i\pmod2\quad(0\leq i<6),
\]
où
\[
 p_5=\frac{s_C-\sum_{i=0}^{4}M_i}{6},\qquad p_i=p_5+M_i\quad(i<5).
\]
Dans ce cas, sa préimage est unique et est donnée par
\[
 e_i=\frac{p_i+a_i}{2},\qquad e_{i+6}=\frac{p_i-a_i}{2}.
\]
\end{lemma}

\begin{proof}
La somme des cinq marges est \(s_C-6p_5\), d'où l'on obtient la
division par six. Une fois les \(p_i\) reconstruits, chaque paire d'équations
\(p_i=e_i+e_{i+6}\), \(a_i=e_i-e_{i+6}\) possède la solution indiquée.
Cette solution est entière exactement sous la congruence de parité. Toutes
les opérations s'inversent, de sorte qu'aucune liberté résiduelle ne subsiste.
\end{proof}

Cette réversibilité montre quelle information est conservée ; elle n’autorise pas à
supprimer le reste de l’article. Le registre complet a un domaine plus grand
que cette carte de recensements. L’extraction des canaux qui alimentent
\(U\) utilise ce registre complet et ne se déduit pas de la liste
\((e_i)\) par une identification de noms.


% BEGIN OWNER /Users/ruben/Documents/New project/output/EXTREMA_MEDIA_RAZON_RECIPROCIDAD_ES_EN_20260918/ARTICULO/ES/source/sections/registro_incidencias.tex
% Suplemento de trabajo. No sustituye registro_k.tex ni cierra por decreto E108.
% Procedencia y alcance: INFORME_RECUPERACION.md, misma carpeta.
\subsection{Évaluation incidencielle antérieure au registre dodécaphasique}
\label{x_reciprocidad:subsec:registro-evaluacion-incidencial}

La récupération d’un registre à partir de ses différences cycliques et de sa
charge totale est une question distincte de l’évaluation de ces coordonnées
sur le registre des incidences. Le premier problème est résolu au moyen de
l’extracteur transversal de la proposition~\ref{x_reciprocidad:prop:k-transversal}.
Pour le second, une présentation historique du
registre fournit trois dénombrements entiers dans chaque fenêtre. Il est utile
d’expliciter leur composition et de déterminer exactement l’information que
requiert leur évaluation.

Soit \(\mathscr L\) un livre de visites et de traces orientées préalablement
généré. Chaque visite conserve la route, la position APP, le feuillet,
l'instant, la multiplicité et l'incidence de frontière. L'évaluation
arithmétique produit l'entier et sa décomposition
\[
 n_t=r_t+9q_t,\qquad 1\leq r_t\leq9.
\]
L'entier \(n_t\) se calcule à partir du feuillet de somme ou de produit ; son
résidu n'est pas reçu comme une valeur indépendante. La distinction entre
incidence de couronne et incidence intérieure demeure explicite pour les
visites de résidu \(9\).

Dans la fenêtre \(E_m\), la présentation par dénombrements considère
\begin{align}
 A_m&=\sum_{v\in E_m}w(v)
             \mathbf1_{\{1,3,5,7\}}(r(v)),\\
 C_m&=\sum_{j\geq0}\bigl(\nu^-_{120,m}(j)-\nu^+_{120,m}(j)\bigr),\\
 V_m&=\sum_{v\in E_m}w(v)\mathbf1_{\{9\}}(r(v))\epsilon_\partial(v).
\end{align}
Ici, \(w(v)\) est la multiplicité d'incidence et \(\epsilon_\partial(v)=1\) dans la couronne, \(-1\) à l'intérieur. Les traces comptées par \(\nu^\pm\) doivent être énumérées au moyen de la règle correspondante de l'article. La somme entière requiert un support fini ou une construction alternative qui détermine son sens. Le nombre d'incidences chronologiques et la longueur réduite \(120(1+3j)\) ont des types distincts ; leur relation exige l'unité de longueur du lecteur.

On définit \([x]_{10}\in\{0,\ldots,9\}\) comme le représentant
euclidien, et l'on conserve les quotients
\[
 x=[x]_{10}+10\lfloor x/10\rfloor.
\]
Le bloc incidenciel est
\begin{equation}
 z_m=100[A_m]_{10}+10[C_m]_{10}+[V_m]_{10}.
 \label{x_reciprocidad:eq:registro-bloque-incidencial}
\end{equation}
Par composition, on obtient l'application
\begin{equation}
 \mathcal E_{\rm cnt}(\mathscr L)=
 \left((z_m-z_{m+3})_{m=1}^{12},
       (z_m-z_{m+4})_{m=1}^{12},\sum_{m=1}^{12}z_m\right).
\label{x_reciprocidad:eq:registro-composicion-incidencial}
\end{equation}
Les indices se lisent cycliquement modulo \(12\). Dans ce paragraphe,
\((D_jz)_m=z_m-z_{m+j}\), pour \(j=3,4\), et
\(Q(z)=\sum_{m=1}^{12}z_m\) ; la composition précédente est donc
\((D_3z,D_4z,Q(z))\). Cette convention est conservée à la
\S\ref{x_reciprocidad:subsec:k-transversal}, où la reconstruction
entière complète est démontrée.
Cette formule produit ses coordonnées à partir des dénombrements. Elle n'utilise
ni \(K\), ni \(U\), ni \(\alpha\), ni les coordonnées transversales attendues
comme arguments. Son identification avec
\(\mathcal E_{108}^{90,120}\) exige de la composer avec la famille et
les incidences effectives de l'état terminal ; elle ne se déduit pas uniquement
de l'inversibilité du système transversal.

\begin{proposition}[Information arithmétique suffisante et nécessaire]
\label{x_reciprocidad:prop:registro-36-residuos}
Soient \(N,N'\in\mathbb Z^{12\times3}\), ordonnés selon les dénombrements
\((A,C,V)\). Pour l'application définie par
\eqref{x_reciprocidad:eq:registro-bloque-incidencial}--\eqref{x_reciprocidad:eq:registro-composicion-incidencial},
\[
 \mathcal E_{\rm cnt}(N)=\mathcal E_{\rm cnt}(N')
 \quad\Longleftrightarrow\quad
 N-N'\in10\mathbb Z^{36}.
\]
En particulier, les \(36\) résidus sectoriels déterminent exactement
la sortie à \(25\) coordonnées. La conservation des quotients et
de la provenance constitue une information supplémentaire du registre.
\end{proposition}

\begin{proof}
L'égalité modulo \(10\) produit les mêmes blocs et, par conséquent,
les mêmes coordonnées transversales. Réciproquement, si les coordonnées
transversales coïncident, \(D_3(z-z')=D_4(z-z')=0\). Les pas
\(3\) et \(4\) engendrent le cycle de longueur \(12\), de sorte que
\(z-z'\) est constant. L'égalité des charges totales annule cette
constante. Enfin, la représentation décimale d'un entier compris entre
\(0\) et \(999\) possède trois chiffres uniques. On récupère ainsi les
trois résidus de chaque fenêtre. L'assertion est une identité algébrique
pour tous \(N,N'\), et non une inférence obtenue par des essais finis.
\end{proof}

\subsection{Conjugaison des dénombrements et termes de frontière}
\label{x_reciprocidad:subsec:registro-conjugacion-incidencial}

Soit \(\rho(m)=13-m\). Considérons une conjugaison d'incidences
qui inverse l'ordre des fenêtres, conserve la classe arithmétique et la
classe de frontière des visites et échange les canaux des traces.
Alors
\[
 (A_m^\dagger,C_m^\dagger,V_m^\dagger)
 =(A_{\rho(m)},-C_{\rho(m)},V_{\rho(m)}).
\]
Si \(c_m=[C_m]_{10}\), la réduction décimale produit l'identité
\begin{equation}
 z_m^\dagger=z_{\rho(m)}
       +100\mathbf1_{\{c_{\rho(m)}\ne0\}}-20c_{\rho(m)}.
 \label{x_reciprocidad:eq:registro-conjugacion-decimal}
\end{equation}
En effet, \([-C]_{10}=0\) lorsque \([C]_{10}=0\), et
\([-C]_{10}=10-[C]_{10}\) dans le cas contraire. La formule conserve
exactement cette différence. La négation du canal n'équivaut pas à nier
le bloc décimal complet.

L'application de cette identité à une involution concrète du TPK exige
de vérifier son action sur les incidences. Pour un lecteur d'occupations
évalué avant l'avancement, l'inversion d'une route
\(\gamma=(x_0,\ldots,x_n)\) satisfait
\[
 \sum_{k=0}^{n-1}f(x_{n-k})
 =\sum_{k=0}^{n-1}f(x_k)+f(x_n)-f(x_0).
\]
Les termes aux extrémités sont incorporés avant la réduction modulo
\(10\). Ils s'annulent sur les routes fermées ; dans les fenêtres ouvertes, ils peuvent
modifier leurs résidus. Cette correction, déjà présente dans le développement
bidirectionnel du transport, spécifie ce que doit conserver le raccordement
entre la route et les dénombrements.

\begin{proposition}[Nécessité de l'information non périodique du registre]
Supposons que l'observable d'une famille fixe de routes satisfasse
\(o(t+54)=o(t)\), que la multiplicité se conserve sous ce retour
et que le même lecteur local agisse dans les fenêtres de \(9\) instants.
Alors ses dénombrements et ses blocs satisfont
\[
 (A,C,V)_{m+6}=(A,C,V)_m,\qquad z_{m+6}=z_m.
\]
\end{proposition}

\begin{proof}
La translation \(t\mapsto t+54\) envoie \(E_m\) bijectivement
sur \(E_{m+6}\), en conservant chaque incidence comptée et sa multiplicité.
L'égalité se transmet à la somme et à la réduction décimale.
\end{proof}

Une présentation à \(12\) blocs distincts exige donc que
le lecteur conserve une information qui ne se factorise pas par cet observable
de période \(54\) : sélection d'incidences, orientation effective,
extrémités, multiplicité ou mémoire. La proposition délimite une réduction
inadmissible de l'état ; elle n'affirme pas que la dynamique complète possède une
période \(54\) et ne détermine pas à elle seule son lecteur terminal.

% END OWNER


% BEGIN OWNER /Users/ruben/Documents/New project/output/EXTREMA_MEDIA_RAZON_RECIPROCIDAD_ES_EN_20260918/ARTICULO/ES/source/sections/censo_terminal_completo.tex
\subsection{Recensement complet des cylindres et sélection de la frontière terminale}
\label{x_reciprocidad:censo:terminal-completo}

La lecture terminale utilise conjointement deux informations : la compatibilité des cylindres régionaux et un profil d'incidence avec des canaux et des positions étiquetés. Dans cette section, on construit les catalogues complets, on énumère leurs triplets et on démontre l'effet de chaque condition du lecteur. Le point de départ est constitué des sorties régionales déjà produites dans la section~\ref{x_reciprocidad:sec:generacion}, et non de valeurs introduites pour définir le transport APP--TRIT--TPK. La réduction à une matrice appartient à une lecture ultérieure de ces sorties et conserve l'ordre des trois canaux.

Le résultat est un théorème de recensement à structure de départ explicite.
Le profil de Gram, des poids, des distances et des marges est déclaré comme donnée du
lecteur terminal. L'exhaustivité de l'énumération démontre ce que sélectionne
ce profil sur les catalogues fixés ; elle ne démontre pas, à elle seule, que le
profil soit une loi universelle ni qu'il puisse être reconstruit à partir du seul triplet
qui le satisfait finalement. Cette séparation permet d'intégrer la preuve
finie complète sans inverser sa généalogie.

\subsubsection{Des sorties régionales aux catalogues ternaires}

Soient \(x_\pi,x_e,x_\varphi\in[0,1)\) les parties fractionnaires des
trois coordonnées régionales déjà construites. On utilise leurs cylindres
publiés de mille positions décimales,
\[
 J_\chi^{(1000)}=
 \left[\frac{d_\chi}{10^{1000}},
       \frac{d_\chi+1}{10^{1000}}\right),
 \qquad \chi\in\{\pi,e,\varphi\}.
\]
Les entiers \(d_\chi\) sont des lectures finies de sortie. La longueur mille
est suffisante pour les entiers qui suivent, comme on le vérifie par les
deux bornes de chaque intervalle ; ce n'est pas une profondeur supposée infinie.
Définissons
\begin{equation}
 p=30,\quad r=1050,\quad L=p+r=1080,\quad k=171,
 \qquad D=3^L,\quad Q=1000^k,
 \label{x_reciprocidad:censo:parametros}
\end{equation}
et
\[
 P_\chi=\lfloor3^{30}x_\chi\rfloor,\qquad
 T_\chi=\lfloor Qx_\chi\rfloor.
\]
Ici, \(k\) compte des triades décimales et ne désigne pas le registre dodécaphasique \(K\). On a \(1000^{171}\leq3^{1080}<1000^{172}\).

\begin{lemma}[Extraction stable d'un entier]
\label{x_reciprocidad:censo:estabilidad}
Pour \(d,m,h\in\mathbb Z\), avec \(m,h>0\), la valeur de
\(\lfloor mx\rfloor\) est constante en \([d/h,(d+1)/h)\) si et seulement si
\begin{equation}
 \left\lfloor\frac{md}{h}\right\rfloor
 =\left\lfloor\frac{m(d+1)-1}{h}\right\rfloor.
 \label{x_reciprocidad:censo:prueba-estabilidad}
\end{equation}
Lorsqu'elle est satisfaite, chacun des deux membres donne cette valeur.
\end{lemma}
\begin{proof}
Le minimum de \(mx\) est atteint à l'extrémité gauche. Son supremum
est \(m(d+1)/h\) et n'est pas atteint. Le plus grand entier que peut prendre
\(\lfloor mx\rfloor\) est
\(\lceil m(d+1)/h\rceil-1\), égal au second membre car
le numérateur et le dénominateur sont des entiers. L'égalité des extrémités
de l'intervalle de valeurs prouve l'affirmation.
\end{proof}

La preuve de stabilité se vérifie sur les trois canaux pour
\(m=3^{30},Q,3^{1080}\). Le dernier choix sera utilisé uniquement
pour une comparaison ultérieure, pas pour sélectionner le profil. Les
premières trente positions récupérées sont précisément les cinq
étapes de six trits du transport régional :
\begin{center}\small
\begin{tabular}{@{}c l r@{}}\toprule
Canal&Préfixe de trente trits&\(P_\chi\)\\\midrule
\(\pi\)&\texttt{010211012222010211002111110221}&29152671743887\\
\(e\)&\texttt{201101121221102011012222102011}&147887858824447\\
\(\varphi\)&\texttt{121200112202121020010210010200}&127247717616687\\\bottomrule
\end{tabular}
\end{center}
La coïncidence conserve la prolongation déjà construite ; la lecture décimale
n'est pas utilisée pour remplacer rétroactivement ses matrices de transport.

Un canal étant fixé, une queue de \(r\) trits se représente par
\(0\leq R<3^r\). Son cylindre et le cylindre des \(k\) premières
triades sont
\begin{equation}
 I_{\chi,R}=\left[\frac{P_\chi3^r+R}{D},
                       \frac{P_\chi3^r+R+1}{D}\right),\qquad
 J_{\chi,T}=\left[\frac{T_\chi}{Q},\frac{T_\chi+1}{Q}\right).
 \label{x_reciprocidad:censo:dos-cilindros}
\end{equation}
Le catalogue inclut toutes les queues pour lesquelles ces deux cylindres
ont une intersection non vide.

\begin{proposition}[Catalogue par intersection exacte]
\label{x_reciprocidad:censo:catalogos}
Les queues compatibles sont les entiers de l'intervalle \([a_\chi,b_\chi]\),
avec
\begin{align}
 a_\chi&=\max\left(0,
       \left\lfloor\frac{T_\chi D}{Q}\right\rfloor-P_\chi3^r\right),
       \label{x_reciprocidad:censo:extremo-a}\\
 b_\chi&=\min\left(3^r-1,
       \left\lfloor\frac{(T_\chi+1)D-1}{Q}\right\rfloor-P_\chi3^r\right).
       \label{x_reciprocidad:censo:extremo-b}
\end{align}
Pour les cylindres régionaux fixés, les troncatures n'interviennent pas et l'on obtient
\begin{equation}
\begin{array}{c|r|r|r}
 \chi&a_\chi\bmod729&b_\chi\bmod729&b_\chi-a_\chi+1\\\hline
 \pi&67&262&196\\
 e&584&51&197\\
 \varphi&117&312&196
\end{array}
\label{x_reciprocidad:censo:tabla-catalogos}
\end{equation}
\end{proposition}
\begin{proof}
Écrivons \(N=P_\chi3^r+R\). Deux intervalles semi-ouverts de
\eqref{x_reciprocidad:censo:dos-cilindros} se coupent exactement lorsque
\[
 NQ<(T_\chi+1)D,\qquad (N+1)Q>T_\chi D.
\]
La deuxième inégalité est équivalente à  
\(N\geq\lfloor T_\chi D/Q\rfloor\) ; la première, à  
\(N\leq\lfloor((T_\chi+1)D-1)/Q\rfloor\). Soustraire  
\(P_\chi3^r\) et conserver \(0\leq R<3^r\) prouve les formules.  
L'évaluation par divisions entières des \(T_\chi\) extraits par  
\eqref{x_reciprocidad:censo:prueba-estabilidad} donne le tableau. Comme \(r\geq6\),  
\(P_\chi3^r\equiv0\pmod{729}\), de sorte que le reste de la première  
extrémité s'obtient également directement de  
\(\lfloor T_\chi D/Q\rfloor\bmod729\).
\end{proof}

Les six derniers trits d'une queue forment le mot
\(\nu_6(R\bmod729)\), où \(\nu_6\) est l'écriture ternaire
de longueur six, y compris ses zéros initiaux. D'après le tableau, chaque
catalogue possède moins de \(729\) éléments consécutifs, de sorte que
cette réduction y est injective. Ainsi, les trois catalogues de frontière
sont entièrement décrits, sans choisir les triplets finaux :
\begin{equation}
 \begin{split}
 \mathcal B_\pi&=\{\nu_6(67+i):0\leq i<196\},\\
 \mathcal B_e&=\{\nu_6((584+j)\bmod729):0\leq j<197\},\\
 \mathcal B_\varphi&=\{\nu_6(117+\ell):0\leq\ell<196\}.
 \end{split}
 \label{x_reciprocidad:censo:fronteras-explicitas}
\end{equation}
Exiger que l'extrémité gauche de \(I_{\chi,R}\) appartienne à
\(J_{\chi,T}\) serait une autre condition : elle perdrait le premier cylindre
qui coupe la frontière et donnerait \(195,196,195\). Le recensement utilise l'intersection des cylindres complets, conformément à sa définition.

\subsubsection{Profil déclaré et indicateurs de compatibilité}

Pour \(u,v\in\mathbb F_3^6\), soient
\(g(u,v)=\sum u_iv_i\in\mathbb F_3\),
\(n(u)=g(u,u)\), \(w(u)=\#\{i:u_i\ne0\}\) et
\(h(u,v)=\#\{i:u_i\ne v_i\}\). Les deux dernières fonctions prennent
des valeurs entières ; les deux premières, des valeurs dans \(\mathbb F_3\).
Le profil local du lecteur, en ordre \((\pi,e,\varphi)\), est
\begin{equation}
 \begin{split}
 (g_{\pi e},g_{\pi\varphi},g_{e\varphi})&=(2,2,0),\qquad
 (n_\pi,n_e,n_\varphi)=(0,0,0),\\
 (w_\pi,w_e,w_\varphi)&=(3,3,3),\qquad
 (h_{\pi e},h_{\pi\varphi},h_{e\varphi})=(2,5,3).
 \end{split}
 \label{x_reciprocidad:censo:perfil-local}
\end{equation}
Les mots sont ordonnés comme dans \eqref{x_reciprocidad:censo:fronteras-explicitas}.
Nous utiliserons \([E]\in\{0,1\}\) pour l'indicateur d'une condition \(E\).
Pour chaque paire de canaux \(\chi,\psi\), la matrice rectangulaire
\(A_{\chi\psi}\) a l'entrée \([g(u,v)=g_{\chi\psi}]\).
La matrice \(A^{h}_{\chi\psi}\) multiplie cette entrée par
\([h(u,v)=h_{\chi\psi}]\). Soient \(D_\chi^n\) la matrice diagonale
de \([n(u)=0]\), et \(D_\chi^{nw}\) celle de
\([n(u)=0][w(u)=3]\).

\begin{theorem}[Énumération exhaustive par sommes finies]
\label{x_reciprocidad:censo:teorema-conteo}
Le produit des trois catalogues contient \(7\,567\,952\) triplets.
Les conditions successives de \eqref{x_reciprocidad:censo:perfil-local} laissent
\begin{equation}
 7\,567\,952\longrightarrow275\,794\longrightarrow6235
 \longrightarrow3127\longrightarrow331.
 \label{x_reciprocidad:censo:cadena-local}
\end{equation}
Ces quatre comptages sont, respectivement,
\begin{align}
 N_g&=\operatorname{tr}(A_{\pi e}A_{e\varphi}A_{\varphi\pi}),
 \nonumber\\
 N_{gn}&=\operatorname{tr}(D_\pi^n A_{\pi e}
             D_e^n A_{e\varphi}D_\varphi^n A_{\varphi\pi}),
 \nonumber\\
 N_{gnw}&=\operatorname{tr}(D_\pi^{nw} A_{\pi e}
             D_e^{nw} A_{e\varphi}D_\varphi^{nw} A_{\varphi\pi}),
 \nonumber\\
 N_{gnwh}&=\operatorname{tr}(D_\pi^{nw} A^h_{\pi e}
             D_e^{nw} A^h_{e\varphi}D_\varphi^{nw} A^h_{\varphi\pi}).
 \label{x_reciprocidad:censo:trazas}
\end{align}
Tous les produits et traces de cette formule sont calculés dans les entiers,  
non dans le corps de trois éléments.
\end{theorem}
\begin{proof}
La taille du produit est \(196\cdot197\cdot196\). En développant
une trace de \eqref{x_reciprocidad:censo:trazas}, ses indices parcourent exactement un
triplet de mots. Le terme vaut un s'il satisfait toutes les conditions
représentées et zéro sinon. Chaque triplet est compté une fois.
Les matrices diagonales ajoutent les conditions individuelles ; les facteurs
rectangulaires ajoutent les conditions par paires. Cela démontre que les
traces sont précisément les cardinaux annoncés.

Pour expliciter son évaluation, une fois fixés les indices \(i,j\) des deux premiers canaux, on forme les ensembles d'indices \(\ell\) du troisième qui satisfont chaque condition par paires. On intersecte leurs indicateurs, on ajoute les indicateurs individuels pertinents et on additionne leur cardinal. La somme finale parcourt \(0\leq i<196\), \(0\leq j<197\). La substitution des mots de \eqref{x_reciprocidad:censo:fronteras-explicitas} et du profil \eqref{x_reciprocidad:censo:perfil-local} donne, dans cet ordre, \(275794,6235,3127,331\). Le supplément enregistre les quatre contributions de chacun des \(196\) indices \(i\) et les \(331\) triplets complets ; son programme évalue ces sommes avec des indicateurs binaires exacts et vérifie scalairement chaque survivant. Aucun échantillon ni aucune tolérance n'intervient.
\end{proof}

\subsubsection{Incidence étiquetée et charge d'orientation}

Pour un triplet survivant, soit \(B\in\mathbb F_3^{3\times6}\) sa
matrice de lignes ordonnées. Le second profil exige
\begin{equation}
 \operatorname{rank}_{\mathbb F_3}B=3,\qquad
 w_{\mathrm{col}}(B)=(1,1,2,2,3,0),\qquad
 q_{\partial}(B)=(1,2,2,1,2,0),
 \label{x_reciprocidad:censo:perfil-columnas}
\end{equation}
où \(w_{\mathrm{col}}\) compte les éléments non nuls de chaque
colonne et \(q_{\partial}\) additionne ses entrées modulo trois. Il s'agit d'un profil
de colonnes étiquetées ; les permuter sans transporter le profil définit un autre
lecteur. Notons \(\mathcal F\) l'ensemble des \(331\) triplets
du théorème précédent.

\begin{proposition}[Réduction du recensement au couple terminal]
\label{x_reciprocidad:censo:dos-matrices}
Les sommes d'indicateurs sur \(\mathcal F\), en exigeant successivement
le rang, les poids de colonne et les sommes de colonne, sont
\begin{equation}
 \sum_{B\in\mathcal F}[\operatorname{rank}B=3]=331,
 \qquad N_{\mathrm{rango,pesos}}=18,
 \qquad N_{\mathrm{rango,pesos,sumas}}=2.
 \label{x_reciprocidad:censo:conteo-columnas}
\end{equation}
Avec des indices de catalogue commençant à un, les deux triplets finaux sont
\((120,197,157)\) et \((129,197,148)\), et leurs matrices sont
\[
 B_0=\begin{pmatrix}0&2&0&2&2&0\\0&0&1&2&2&0\\1&0&1&0&1&0\end{pmatrix},
 \qquad
 B_1=\begin{pmatrix}0&2&1&0&2&0\\0&0&1&2&2&0\\1&0&0&2&1&0\end{pmatrix}.
\]
\end{proposition}
\begin{proof}
Dans chaque triplet de l'ensemble déjà énuméré, on calcule le rang par élimination en \(\mathbb F_3\), on compte les éléments non nuls des six colonnes et on les additionne modulo trois. Les produits des indicateurs de \eqref{x_reciprocidad:censo:perfil-columnas} donnent \eqref{x_reciprocidad:censo:conteo-columnas}.
La substitution des deux indices finaux en
\eqref{x_reciprocidad:censo:fronteras-explicitas} donne les matrices exposées. La liste
complète des survivants et les filtres exacts du supplément permettent
de reproduire l'élimination sans supposer ces deux matrices comme domaine.
\end{proof}

L'orientation \(\sigma=(1,1,-1)=(1,1,2)\in\mathbb F_3^3\)
sélectionne la droite
\(\langle\sigma\rangle=\{(0,0,0),(1,1,2),(2,2,1)\}\).
Les charges des lignes de \(B_0,B_1\) sont \((0,2,0)\) et \((2,2,1)\) ;
seule la seconde appartient à cette droite. On retrouve ainsi la dernière étape
du sélecteur de la section~\ref{x_reciprocidad:subsec:k-terminal}, désormais précédée de
la construction et de l'énumération de son domaine complet. La lecture stable
de \(\lfloor3^{1080}x_\chi\rfloor\), effectuée après les filtres,
donne les frontières \texttt{021020}, \texttt{001220}, \texttt{100210},
avec les mêmes rangs \((129,197,148)\). Cette comparaison ultérieure ne définit
pas le profil et ne participe pas aux sommes d'indicatrices.

\subsubsection{Information conservée et portée du résultat}

Le procédé conserve le préfixe de trente trits, la queue compatible, l'ordre régional, la position de chaque trit et l'orientation. La réduction modulo \(729\) conserve ici l'identité de chaque candidat car l'intervalle des queues a une longueur inférieure à \(729\) ; cette injectivité n'est pas affirmée pour une histoire arbitraire. Le recensement de \(331\) démontre également que le profil local isolé ne sélectionne pas une seule matrice : les marges étiquetées et la charge apportent des conditions effectives supplémentaires.

On a donc récupéré une énumération finie intégrale et son sélecteur sur des cylindres de sortie, avec un profil de lecture explicite. La construction du registre signé utilise en outre son lecteur de mémoire et ses canaux, comme cela est spécifié ci-dessous ; la matrice visible ne se confond pas avec ce registre par égalité du nombre de composants. Ce compte ne démontre pas non plus l'universalité du profil, ne génère pas à nouveau les coordonnées régionales et ne décide pas d'une affirmation sur alpha ou sur le continu complet.

% END OWNER


\subsection{Deux lectures de l'état terminal}
\label{x_reciprocidad:subsec:k-terminal}

Désignons par \(x_{\mathrm{term}}\) l'état produit par la composition de sélection, de transport et de mise à jour jusqu'à la coupe terminale. L'écriture opératorielle est
\[
 x_{\mathrm{term}}=
 (\mathcal U_{108}\circ\cdots\circ\mathcal U_1)(x_{\mathrm{seed}}),
 \qquad \mathcal U_t=\mathsf{Upd}_t\circ\mathsf{Tra}_t\circ\mathsf{Sel}_t.
\]
Ses deux lectures pertinentes sont
\[
 x_{\mathrm{term}}
 \xrightarrow{(\pi_{\mathrm{term}},R_{\mathrm{sgn}})}
 (B_{\mathrm{term}},U).
\]
La matrice \(B_{\mathrm{term}}\) retient la lecture ternaire visible ; le vecteur \(U\) conserve une autre coordonnée de la même histoire. On n'insère pas entre eux une flèche \(B_{\mathrm{term}}\to U\).

Le sélecteur terminal opère sur les trois catalogues régionaux
de tailles 196, 197 et 196, construits dans la
section~\ref{x_reciprocidad:censo:terminal-completo}. Les signatures locales réduisent leurs
\(196\cdot197\cdot196=7\,567\,952\) triplets à 331, et la marge des
colonnes \(q_{\partial}=(1,2,2,1,2,0)\) laisse deux matrices :
\[
 B_0=\begin{pmatrix}0&2&0&2&2&0\\0&0&1&2&2&0\\1&0&1&0&1&0\end{pmatrix},
 \qquad
 B_1=\begin{pmatrix}0&2&1&0&2&0\\0&0&1&2&2&0\\1&0&0&2&1&0\end{pmatrix}.
\]
L'orientation des lignes est \(\sigma=(1,1,-1)=(1,1,2)\) dans
\(\mathbb F_3^3\). Les charges des lignes des candidates sont
\((0,2,0)\) et \((2,2,1)\). Seule la seconde appartient à
\(\langle\sigma\rangle=\{(0,0,0),(1,1,2),(2,2,1)\}\) ; par conséquent,
\(B_{\mathrm{term}}=B_1\). Cette vérification prouve la dernière étape du
sélecteur à partir du recensement déclaré. Elle ne transforme pas la marge des colonnes
en une conséquence de la charge des lignes, et ne remplace pas l'énumération
précédente par l'inspection de ces deux matrices.

Le lecteur signé possède la composition typée
\begin{equation}
 R_{\mathrm{sgn}}
 =\Pi_H\circ\mathcal E_{108}^{90,120}\circ\mathcal R_{12}:
 \mathcal X_{\mathrm{TPK}}^{\mathrm{term,mem}}
 \longrightarrow\mathbb Z^{12}.
 \label{x_reciprocidad:eq:k-lector-firmado}
\end{equation}
Le registre des canaux détermine
\[
 \mathcal E_{108}^{90,120}(\mathcal R_{12}(x_{\mathrm{term}}))
 =(b^{90},b^{120},Q_{\mathrm{TPK}}),
\]
avec les coordonnées
\begin{align*}
 b^{90}={}&(-495,-116,-684,108,601,-90,\\
            &\hspace{19mm}-173,-88,313,560,-397,461),\\
 b^{120}={}&(-425,-281,-481,671,-255,30,\\
             &\hspace{19mm}475,-543,680,251,6,-128),\\
 Q_{\mathrm{TPK}}={}&6263.
\end{align*}
Ces quantités sont utilisées comme coordonnées du registre préalable, non
comme paramètres ajustés au moyen d'alpha. La reconstruction démontrée
dans les paragraphes suivants commence exactement à cette sortie typée.
La caractérisation intégrale de l'extracteur, son inverse et la composition
des contributions sont développées dans la section~\ref{x_reciprocidad:img09:seccion},
en particulier dans le théorème~\ref{x_reciprocidad:img09:imagen} et la
proposition~\ref{x_reciprocidad:img09:acumulacion}. La détermination régionale
du descripteur et la sélection terminale sont exposées dans le
chapitre~\ref{chap:despliegue-selector-k}. Chaque construction conserve son
domaine : le sélecteur agit sur les répertoires engendrés et l'inversion
agit sur les canaux signés déclarés ici. La trace documentaire de
l'extracteur est consignée dans la provenance technique du paquet.

\subsection{Construction du registre signé au moyen des canaux orientés}
\label{x_reciprocidad:subsec:k-pi-h}

Le lecteur harmonique \(\Pi_H\) peut s'écrire sans utiliser ni \(K\) ni
alpha. Pour éviter de confondre ses sommes orbitales avec les blocs d'alpha,
nous les noterons \(s_1,s_2,s_3\). Soient
\[
 B_r=\sum_{j=0}^{3}b^{120}_{r+3j},\qquad
 d_{r,j}=b^{90}_{r+3j},\qquad1\leq r\leq3,
\]
et
\begin{equation}
 s_1=\frac{Q_{\mathrm{TPK}}+2B_1+B_2}{3},\qquad
 s_2=s_1-B_1,\qquad s_3=s_2-B_2.
 \label{x_reciprocidad:eq:k-sumas-armonicas}
\end{equation}
La division par trois appartient au lecteur des trois orbites. Son caractère entier est vérifié sur les sorties compatibles du registre ; il n'est pas supposé pour un élément arbitraire de \(\mathbb Z^{25}\).

Les trois blocs signés sont
\begin{equation}
 U^{(r)}=(s_r,d_{r,1}-d_{r,3},d_{r,0}+d_{r,2},d_{r,0}-d_{r,2}).
 \label{x_reciprocidad:eq:k-bloques-firmados}
\end{equation}
La substitution entière donne
\[
 (B_1,B_2,B_3)=(972,-1073,101),\qquad
 (s_1,s_2,s_3)=(2378,1406,2479),
\]
et, par conséquent,
\begin{align*}
 U^{(1)}&=(2378,-452,-668,-322),\\
 U^{(2)}&=(1406,998,-204,-28),\\
 U^{(3)}&=(2479,-551,-371,-997).
\end{align*}
En entrelaçant les colonnes, on obtient
\begin{equation}
 U=(2378,1406,2479,-452,998,-551,
       -668,-204,-371,-322,-28,-997).
 \label{x_reciprocidad:eq:k-u-firmado}
\end{equation}
La présence de valeurs négatives et de composantes supérieures à 999
montre que \(U\) n'est pas un mot de chiffres en base mille. Ce n'est pas non plus
\(U_6\Vert U_6\), concaténation du catalogue de période six. La différence
concerne le domaine et l'information conservée, et non seulement les valeurs.

L'ordre constructif jusqu'ici est
\[
 x_{\mathrm{term}}\longrightarrow\mathcal R_{12}
 \longrightarrow(b^{90},b^{120},Q_{\mathrm{TPK}})
 \longrightarrow U.
\]
L'inversion qui suit produit le registre dodécaphasique. Le fait que nous puissions ensuite récupérer les canaux à partir du registre dodécaphasique constitue un contrôle de réversibilité ; il n'inverse pas cet ordre causal.

\subsection{Hadamard par orbites et reconstruction entière}
\label{x_reciprocidad:subsec:k-hadamard}

Dans le cycle à douze positions, le pas trois détermine trois orbites :
\[
 (1,4,7,10),\qquad(2,5,8,11),\qquad(3,6,9,12).
\]
La séparation de quatre secteurs de chaque orbite est celle qui, une fois
l'origine et l'orientation fixées, admet la lecture angulaire de quatre-vingt-dix degrés. La
transformation n'agit pas sur trois quadruplets contigus de l'ordre naturel.

Soit
\[
 H_4=\begin{pmatrix}
 1&1&1&1\\1&1&-1&-1\\1&-1&1&-1\\1&-1&-1&1
 \end{pmatrix}.
\]
Si \(P_{\mathrm{orb}}\) réordonne selon les trois orbites indiquées, nous définissons
\[
 \mathcal H_{12}=P_{\mathrm{orb}}^{-1}
 (H_4\oplus H_4\oplus H_4)P_{\mathrm{orb}}.
\]

\begin{lemma}[Inversion de Hadamard et intégralité]
\label{x_reciprocidad:lem:k-hadamard}
On a
\[
 H_4^2=4I_4,\qquad H_4^{-1}=\frac14H_4,\qquad\det H_4=-16.
\]
Pour \(u\in\mathbb Z^4\), la préimage \(H_4^{-1}u\) est entière si et
seulement si \(H_4u\equiv0\pmod4\).
\end{lemma}

\begin{proof}
Les lignes sont orthogonales et le carré de leur norme vaut quatre ; la matrice est symétrique. Cela donne \(H_4^2=4I_4\) et l'inverse. Une élimination entière, ou le développement du déterminant, donne le signe négatif et la valeur \(-16\). La condition de caractère entier est exactement la divisibilité par quatre des quatre coordonnées de \(H_4u\).
\end{proof}

\begin{definition}[Transformation entière dodécaphasique]
Le sous-réseau des registres signés et sa coordinatisation entière sont
\[
 \mathcal L_H=\{u\in\ZZ^{12}:\mathcal H_{12}u\equiv0\pmod4\},
 \qquad \mathscr K(u)=\tfrac14\mathcal H_{12}u.
\]
L'application \(\mathscr K:\mathcal L_H\to\ZZ^{12}\) est un isomorphisme
de groupes abéliens, d'inverse \(k\mapsto\mathcal H_{12}k\).
Le registre canonique \(K\) est l'image du registre signé généré
par les opérations précédentes, avec une origine de phase et une orientation fixées.
\end{definition}

\begin{proof}
La congruence définit exactement l'intégralité de \(\mathscr K(u)\).
L'identité \(\mathcal H_{12}^2=4I\) démontre les deux compositions
inverses. En particulier, tout \(k\in\ZZ^{12}\) a pour préimage
\(\mathcal H_{12}k\in\mathcal L_H\).
\end{proof}

\begin{proposition}[Registre entier canonique]
\label{x_reciprocidad:prop:k-sello}
L'inversion \(K^{(r)}=H_4U^{(r)}/4\) des blocs de
\eqref{x_reciprocidad:eq:k-bloques-firmados} produit
\begin{align*}
 K^{(1)}&=(234,729,621,794),\\
 K^{(2)}&=(543,659,58,146),\\
 K^{(3)}&=(140,824,914,601).
\end{align*}
Dans l'ordre naturel, le registre dodécaphasique est
\begin{equation}
 K=(234,543,140,729,659,824,621,058,914,794,146,601).
 \label{x_reciprocidad:eq:k-vector}
\end{equation}
C'est l'unique préimage rationnelle de \(U\), et elle appartient à \(\{0,\ldots,999\}^{12}\).
\end{proposition}

\begin{proof}
Les trois produits avant division sont
\begin{align*}
 H_4U^{(1)}&=(936,2916,2484,3176),\\
 H_4U^{(2)}&=(2172,2636,232,584),\\
 H_4U^{(3)}&=(560,3296,3656,2404).
\end{align*}
Toutes les coordonnées sont divisibles par quatre. Diviser et défaire la
permutation des orbites donne \eqref{x_reciprocidad:eq:k-vector}. L'unicité procède de l'
inversibilité rationnelle, et non d'une recherche parmi des candidats proches d'
une valeur physique.
\end{proof}

Le caractère entier peut également se formuler en termes de réseaux. Les diviseurs déterminantaux de \(H_4\) sont \(1,2,4,16\) : les entrées ont pour plus grand commun diviseur un ; les mineurs d'ordre deux sont pairs et l'un d'eux est de module deux ; ceux d'ordre trois sont des multiples de quatre et l'un d'eux est de module quatre. On en déduit
\[
 \operatorname{SNF}(H_4)=\operatorname{diag}(1,2,2,4),
\]
et pour la transformation globale,
\[
 \operatorname{coker}\mathcal H_{12}
 \cong(\mathbb Z/2\mathbb Z)^6\oplus(\mathbb Z/4\mathbb Z)^3.
\]
L'indice de l'image est \(16^3=4096\). Ainsi, tout vecteur signé
entier n'admet pas nécessairement une préimage entière ; le vecteur obtenu satisfait les
congruences précises exigées par l'inversion. Le facteur \(1/4\) est
imposé par la matrice et n'a pas le statut d'un coefficient libre.

\subsection{Différences cycliques et composante uniforme}
\label{x_reciprocidad:subsec:k-transversal}

Avec des indices cycliques modulo douze, soient
\[
 (D_3z)_m=z_m-z_{m+3},\qquad(D_4z)_m=z_m-z_{m+4},
 \qquad Q(z)=\sum_{m=1}^{12}z_m.
\]
Le premier opérateur compare des secteurs séparés par trois positions ; le
second, par quatre. Les noms de canal 90 et 120 se réfèrent ici à
ces séparations dans le cycle orienté, et non à une identification avec un
opérateur physique d'un autre domaine.

\begin{proposition}[Complétude du système transversal]
\label{x_reciprocidad:prop:k-transversal}
L'opérateur
\[
 \mathcal D=(D_3,D_4,Q):\mathbb Z^{12}\longrightarrow\mathbb Z^{25}
\]
a un rang égal à douze et des facteurs invariants non nuls
\[
 \operatorname{diag}(1,\ldots,1,12),
\]
avec onze unités. En particulier, leurs sorties déterminent un unique registre dodécaphasique.
\end{proposition}

\begin{proof}
Si \(D_3z=D_4z=0\), le vecteur est constant le long des deux pas.
Comme \(3\) et \(4\) engendrent le cycle modulo douze,
\(\ker D_3\cap\ker D_4=\langle\mathbf1\rangle\). La charge totale
élimine cette dernière liberté parce que \(Q(\mathbf1)=12\).

Pour l'affirmation entière, les lignes de \((D_3,D_4)\) sont des incidences orientées d'un graphe connexe. Un arbre couvrant fournit onze facteurs invariants unitaires. Tout mineur d'ordre douze doit contenir la ligne \(Q\). En ajoutant les onze premières colonnes à la dernière, celle-ci s'annule dans les lignes d'incidence et vaut douze dans la ligne \(Q\). Tous les mineurs maximaux sont divisibles par douze et l'arbre en produit un de module exactement douze. Le dernier facteur est donc douze.
\end{proof}

La vérification sur \eqref{x_reciprocidad:eq:k-vector} donne
\[
 D_3K=b^{90},\qquad D_4K=b^{120},\qquad Q(K)=6263.
\]
Les formules \eqref{x_reciprocidad:eq:k-sumas-armonicas} et
\eqref{x_reciprocidad:eq:k-bloques-firmados} reconstruisent à leur tour \(U\) à partir de ces
différences. On ferme ainsi un triangle de représentations exactes du
même objet : registre transversal, vecteur signé et registre dodécaphasique. L'égalité
des résultats ne fait pas de leurs codomaines un seul et même domaine ; elle démontre que
les applications déclarées restituent l'information pertinente.


% BEGIN OWNER /Users/ruben/Documents/New project/output/EXTREMA_MEDIA_RAZON_RECIPROCIDAD_ES_EN_20260918/ARTICULO/ES/source/sections/registro_imagen_integral.tex
% Formalización nueva integrada en la revisión III; originales preservados.
% Propietarios: U016, ecuaciones extractor-diferencias-k y lector-armonico-*;
% registro_k.tex, subsecciones k-pi-h, k-hadamard y k-transversal.
\subsection{Compatibilité entière et détermination du registre transversal}
\label{x_reciprocidad:img09:seccion}

Le domaine de cette construction est la publication du registre d'une
histoire APP--TRIT--TPK. On distingue l'évaluation des événements du
registre, la compatibilité de ses canaux et l'inversion de ses
coordonnées. Le résultat suivant donne un critère entier complet pour
le deuxième niveau et une factorisation explicite pour y relier le premier.
La démonstration utilise les opérateurs du registre et établit leurs
relations pour tout élément des domaines entiers indiqués.

\subsubsection{Caractérisation par accroissements adjacents}

Les indices appartiennent à $\mathbb Z/12\mathbb Z$. Soient
\[
 (Sz)_i=z_{i+1},\qquad D_3=I-S^3,\qquad D_4=I-S^4,
 \qquad Q(z)=\sum_{i=0}^{11}z_i.
\]
On conserve l'extracteur transversal
\[
 \mathcal D:\mathbb Z^{12}\longrightarrow\mathbb Z^{25},
 \qquad k\longmapsto(D_3k,D_4k,Q(k)).
\]
Pour un triplet de coordonnées entières $(b,c,q)$, on écrit
\begin{equation}
 w_i=c_i-b_{i+1},\qquad
 P_0=0,\quad P_i=\sum_{j=0}^{i-1}w_j\ (1\leq i\leq12),
 \qquad T(w)=\sum_{i=0}^{11}P_i.
 \label{x_reciprocidad:img09:incremento}
\end{equation}

\begin{theorem}[Image entière de l'extracteur transversal]
\label{x_reciprocidad:img09:imagen}
Un triplet $(b,c,q)\in\mathbb Z^{25}$ appartient à
$\operatorname{im}\mathcal D$ si et seulement s'il satisfait
\begin{align}
 b_i&=w_i+w_{i+1}+w_{i+2}&& (0\leq i<12),
 \label{x_reciprocidad:img09:relaciones}\\
 \sum_{i=0}^{11}w_i&=0,
 \label{x_reciprocidad:img09:cierre}\\
 q+T(w)&\equiv0\pmod {12}.
 \label{x_reciprocidad:img09:congruencia}
\end{align}
Son unique préimage entière est
\begin{equation}
 k_i=\frac{q+T(w)}{12}-P_i,\qquad0\leq i<12.
 \label{x_reciprocidad:img09:inversa}
\end{equation}
Les douze équations de \eqref{x_reciprocidad:img09:relaciones} et l'équation
\eqref{x_reciprocidad:img09:cierre} sont linéairement indépendantes sur $\mathbb Q$.
\end{theorem}
\begin{proof}
Pour une sortie de $\mathcal D$, on a
\[
 c_i-b_{i+1}=(k_i-k_{i+4})-(k_{i+1}-k_{i+4})=k_i-k_{i+1}.
\]
La somme de trois accroissements donne $b_i$ ; la somme cyclique de tous
est nulle. De plus, $P_i=k_0-k_i$, d'où
$q=12k_0-T(w)$. On obtient toutes les conditions et la formule inverse.

Réciproquement, \eqref{x_reciprocidad:img09:congruencia} rend entière la formule
\eqref{x_reciprocidad:img09:inversa} ; \eqref{x_reciprocidad:img09:cierre} permet de la prolonger
cycliquement. Ses différences adjacentes sont $w_i$. Par conséquent,
$D_3k=b$ et
\[
 (D_4k)_i=w_i+w_{i+1}+w_{i+2}+w_{i+3}
 =w_i+b_{i+1}=c_i.
\]
La somme de l'inverse donne $Q(k)=q$. La même formule prouve l'unicité.
Pour l'indépendance, le changement entier inversible
$(b,c)\mapsto(b,w=c-Sb)$ transforme les douze premières relations en
$b=(I+S+S^2)w$. Chacune isole une coordonnée différente de $b$.
La condition $\sum w_i=0$ est une relation supplémentaire indépendante.
\end{proof}

L'image rationnelle a pour dimension douze. Sa structure entière ajoute
une congruence d'ordre douze : celle-ci est la manifestation constructive du
dernier facteur de Smith de $\mathcal D$. Les onze accroissements
$w_0,\ldots,w_{10}$ et la charge $q$ sont des coordonnées suffisantes, avec
$w_{11}=-\sum_{i=0}^{10}w_i$ et la congruence
\eqref{x_reciprocidad:img09:congruencia}. Les vingt-cinq coordonnées initiales
conservent ces mêmes douze degrés de liberté au moyen de treize relations.

\subsubsection{Commutativité avec la lecture harmonique}

Soit $\mathcal H_{12}$ la transformation de Hadamard du registre,
avec les orbites $(r,r+3,r+6,r+9)$ pour $r=0,1,2$. Son bloc est
\[
 H_4=\begin{pmatrix}
 1&1&1&1\\1&1&-1&-1\\1&-1&1&-1\\1&-1&-1&1
 \end{pmatrix},\qquad H_4^2=4I_4.
\]
Pour $(b,c,q)$ compatible, on définit
\[
 B_r=\sum_{j=0}^3c_{r+3j},\qquad
 A_0=\frac{q+2B_0+B_1}{3},\quad A_1=A_0-B_0,
 \quad A_2=A_1-B_1.
\]
Le lecteur $\Pi_H$ de \eqref{x_reciprocidad:eq:k-bloques-firmados} possède les blocs
\[
 (\Pi_H(b,c,q))^{(r)}=
 (A_r,b_{r+3}-b_{r+9},b_r+b_{r+6},b_r-b_{r+6}).
\]

\begin{proposition}[Concordance des deux reconstructions]
\label{x_reciprocidad:img09:triangulo}
Sur $\operatorname{im}\mathcal D$, on a
\[
 \Pi_H\mathcal D=\mathcal H_{12},\qquad
 \mathcal D^{-1}=\tfrac14\mathcal H_{12}\Pi_H.
\]
L'image de $\Pi_H$ restreinte à ce domaine est exactement
\[
 \mathcal L_H=
 \{u\in\mathbb Z^{12}:\mathcal H_{12}u\equiv0\pmod4\}.
\]
\end{proposition}
\begin{proof}
Pour $k=\mathcal D^{-1}(b,c,q)$, soit
$a_r=\sum_{j=0}^3k_{r+3j}$. Le décalage de quatre positions
envoie chaque orbite sur la suivante, et donc $B_r=a_r-a_{r+1}$.
Avec $q=a_0+a_1+a_2$, ces identités donnent $A_r=a_r$.
Sur une orbite, notons $(x_0,x_1,x_2,x_3)$ les coordonnées
de $k$ et $d_j=x_j-x_{j+1}$. Alors
\[
 d_1-d_3=x_0+x_1-x_2-x_3,\quad
 d_0+d_2=x_0-x_1+x_2-x_3,\quad
 d_0-d_2=x_0-x_1-x_2+x_3.
\]
Ce sont les trois autres lignes de $H_4k^{(r)}$. L'inversion et la
caractérisation entière découlent de $\mathcal H_{12}^2=4I$.
\end{proof}

La condition $\Pi_H(b,c,q)\in\mathcal L_H$ ne contrôle que la
sortie harmonique. L'appartenance des vingt-cinq coordonnées à
$\operatorname{im}\mathcal D$ exige en outre
\eqref{x_reciprocidad:img09:relaciones}--\eqref{x_reciprocidad:img09:congruencia}.
Par exemple, ajouter à $c$ le vecteur $e_0-e_3$ conserve les trois sommes
orbitales et laisse $\Pi_H$ inchangé. En partant de $(0,0,0)$, on
obtient ainsi un triplet avec $\Pi_H=0$ qui ne satisfait pas
\eqref{x_reciprocidad:img09:relaciones}. Il s'agit d'un contrôle négatif exact de la
compatibilité du canal, indépendant de toute évaluation terminale.

\subsubsection{Détermination et transport sur les orbites du registre}

Le théorème~\ref{x_reciprocidad:img09:imagen} détermine $k$ lorsque ses
canaux et sa charge ont été produits. Comme condition sur des sorties non encore
individualisées, la même image est un groupe isomorphe à
$\mathbb Z^{12}$. Une fois fixée uniquement une charge $q$, les solutions
forment une classe affine de
\[
 L_0=\{v\in\mathbb Z^{12}:\sum v_i=0\}
 =\bigoplus_{i=0}^{10}\mathbb Z(e_i-e_{11}).
\]
Par exemple, $100\mathbf1+e_0$ et $100\mathbf1+e_1$ sont deux
registres distincts, de charge $1201$, avec des composantes dans
$\{0,\ldots,999\}$. Chacun satisfait toutes les conditions
d'image et toutes les congruences de Hadamard au moyen de ses propres
canaux. Ces témoins comparent les équations de compatibilité ;
on ne leur attribue pas le statut d'histoires terminales HMT.

\begin{proposition}[Covariance cyclique et stabilisateur du registre]
\label{x_reciprocidad:img09:covariancia}
Soient $\rho$ la translation des fenêtres sur un domaine de registres
$\mathscr R$ et $S$ la translation précédente sur $\mathbb Z^{12}$.
Sur l'orbite de $R\in\mathscr R$, un lecteur covariant est
déterminé par un vecteur
\[
 k\in\operatorname{Fix}(S^p),
\]
où $p\mid12$ est la période minimale de $R$ sous $\rho$.
La condition de charge $Q(k)=q$ admet des solutions entières si et seulement si
$12/p$ divise $q$ ; lorsqu'elles existent, elles forment une classe affine de rang
$p-1$. En particulier, sur une orbite libre, covariance et charge
laissent onze degrés de liberté.
\end{proposition}
\begin{proof}
La règle $F(\rho^jR)=S^jk$ est bien définie exactement lorsque
$S^pk=k$. Ces vecteurs répètent $12/p$ fois un mot de
longueur $p$. Leur charge est $\frac{12}{p}\sum_{i=0}^{p-1}k_i$.
La divisibilité et le rang découlent de cette formule. La
covariance des canaux résulte de $D_aS=SD_a$ et de $QS=Q$.
Sur $\mathcal L_H$, l'action correspondante est
$\widehat S_H=\mathcal H_{12}S\mathcal H_{12}/4$ ;
d'après la proposition~\ref{x_reciprocidad:img09:triangulo}, les deux reconstructions
entrelacent ces actions.
\end{proof}

Si le domaine porte une origine marquée, il faut utiliser l'action qui
transporte aussi cette marque. La proposition concerne cette action
cyclique déclarée ; elle ne présuppose pas la périodicité de l'histoire
enrichie. La naturalité par rapport aux troncatures ultérieures à la coupure
$108$ s'obtient, pour tout lecteur $F_{108}$ déjà défini, au moyen de
$F_N=F_{108}\circ\tau_{N,108}$. Cette naturalité conserve la règle
d'évaluation initiale ; elle ne choisit pas à elle seule l'une des règles possibles
sur les cent huit premiers événements.

\subsubsection{Évaluation par événements et composition des contributions}

La caractérisation précédente fournit une factorisation de l'extraction au moyen de deux
lectures effectives du registre :
\begin{equation}
 \omega:\mathscr R\longrightarrow
 \{w\in\mathbb Z^{12}:\sum w_i=0\},\qquad
 q_{\rm reg}:\mathscr R\longrightarrow\mathbb Z,
 \quad q_{\rm reg}(R)+T(\omega(R))\equiv0\pmod {12}.
 \label{x_reciprocidad:img09:contrato-productor}
\end{equation}
Ces applications doivent être évaluées sur les événements, l'orientation,
l'incidence et la mémoire produits par APP--TRIT--TPK. Une fois données leurs
actions sur ces données, la composition
\[
 R\longmapsto(w,q)\longmapsto
 ((I+S+S^2)w,(I+S+S^2+S^3)w,q)
 \xrightarrow{\Pi_H}u\xrightarrow{\mathcal H_{12}/4}k
\]
est déterminée et satisfait tous les contrôles inverses. La formule
$w=c-Sb$ permet également de raccorder un producteur qui fournit déjà les deux
canaux : c'est une vérification portant sur sa sortie, antérieure à toute comparaison
avec le témoin terminal.

\begin{proposition}[Accumulation des contributions et unicité prospective]
\label{x_reciprocidad:img09:acumulacion}
Supposons que l'état initial et chaque événement $a_t$ du registre
produisent, selon des règles fixées au préalable, des contributions
$(\delta w_t,\delta q_t)$ avec
\[
 \sum_i\delta w_{t,i}=0,\qquad
 \delta q_t+T(\delta w_t)\equiv0\pmod {12}.
\]
Avec des conditions initiales compatibles $(w^{(0)},q^{(0)})$,
l'actualisation
\[
 w^{(t+1)}=w^{(t)}+\delta w_t,\qquad
 q^{(t+1)}=q^{(t)}+\delta q_t
\]
produit un registre unique à chaque étape. Son accroissement est
\[
 \delta k_{t,i}=
 \frac{\delta q_t+T(\delta w_t)}{12}-P_i(\delta w_t).
\]
La construction respecte la concaténation des registres, et la lecture
d'un préfixe reste identique sous les prolongements ultérieurs.
\end{proposition}
\begin{proof}
Les conditions de compatibilité se conservent par addition. La formule
\eqref{x_reciprocidad:img09:inversa} est additive en $(w,q)$ et produit
l'accroissement indiqué. La récurrence fixe successivement chaque état.
La somme sur deux segments est la somme concaténée, et un préfixe reçoit
exactement les mêmes termes avant et après le prolongement de l'histoire.
\end{proof}

La proposition fournit un critère suffisant, et non une exigence de
linéarité pour tout lecteur TPK. Si les événements sont transportés par des actions
non triviales, leurs contributions peuvent s’exprimer dans la carte terminale
au moyen du cocycle affine correspondant. L’équation de cocycle conserve
alors la même indépendance à l’égard de la partition de la trajectoire.
Dans les deux cas, les règles d’événement et l’état initial constituent les
données constructives qui doivent provenir de l’article : les choisir pour
reproduire un vecteur attendu remplacerait cette provenance par une autre tâche.

Le résultat de ce développement est la condition exacte de détermination
\eqref{x_reciprocidad:img09:contrato-productor}, son inversion démontrée et ses identités
de transport. Sa portée est le registre transversal ; les autres
constructions de HMT conservent leurs domaines et leurs preuves propres.

% END OWNER


\subsection{Deux lectures rationnelles : \texorpdfstring{\(\kappa_{36}\)}{kappa36} et \texorpdfstring{\(\kappa_{\mathrm{per}}\)}{kappa périodique}}
\label{x_reciprocidad:subsec:k-kappas}

Une fois \(K\) construit, la carte en base mille permet deux opérations
distinctes. La première lit un seul tour ; la seconde prolonge le registre dodécaphasique
périodiquement. Soit \(B=1000\) et définissons son entier de concaténation
\begin{equation}
 N_K=\sum_{m=1}^{12}K_mB^{12-m}
 =234543140729659824621058914794146601.
 \label{x_reciprocidad:eq:k-numerador}
\end{equation}
Les zéros initiaux d'un bloc, comme celui de 058, font partie de son écriture
à trois chiffres. Ils ne changent pas sa valeur entière, mais permettent de concaténer
sans ambiguïté les douze blocs.

\begin{definition}[Lectures finie et périodique du registre dodécaphasique]
\label{x_reciprocidad:def:k-kappas}
La lecture d'un tour est
\[
 \kappa_{36}=\sum_{m=1}^{12}K_mB^{-m}=\frac{N_K}{B^{12}}.
\]
Le prolongement périodique est
\[
 K_j^{\mathrm{per}}=K_{1+((j-1)\bmod12)},\qquad
 \kappa_{\mathrm{per}}=\sum_{j\geq1}K_j^{\mathrm{per}}B^{-j}.
\]
\end{definition}

\begin{proposition}[Valeur et période exactes du registre dodécaphasique prolongé]
\label{x_reciprocidad:prop:k-periodo}
On a
\begin{equation}
 \kappa_{\mathrm{per}}=\frac{N_K}{B^{12}-1},\qquad
 \kappa_{\mathrm{per}}-\kappa_{36}
 =\frac{N_K}{B^{12}(B^{12}-1)}
 =B^{-12}\kappa_{\mathrm{per}}.
 \label{x_reciprocidad:eq:k-diferencia-kappas}
\end{equation}
Le mot infini du registre dodécaphasique a pour période minimale douze en base mille et pour période minimale trente-six en base dix.
\end{proposition}

\begin{proof}
La somme du premier tour est \(N_KB^{-12}\) ; chaque nouveau tour
multiplie sa contribution par \(B^{-12}\). La série géométrique produit
\(N_K/(B^{12}-1)\), et la soustraction donne la seconde identité.

Les douze composantes de \(K\) sont distinctes. Une période propre du cycle
identifierait deux d'entre elles ; la période minimale des blocs est donc
douze. Le développement décimal est canonique, car le registre dodécaphasique n'a pas de queue
de blocs 999. Si sa période décimale minimale était \(d\), elle diviserait
36. Un regroupement par trois produirait une période de blocs
\(d/\gcd(d,3)\). Aucun diviseur propre de 36 ne donne douze par cette
opération. La période décimale minimale est donc 36.
\end{proof}

Les deux lectures sont rationnelles, mais ne sont pas le même rationnel. La première se termine après douze blocs ; la seconde répète le tour. En outre, \(0<\kappa_{\mathrm{per}}-\kappa_{36}<B^{-12}\). Leur proximité ne permet pas de substituer l'une à l'autre dans une identité exacte. En particulier, la clôture finie d'alpha utilise \(\kappa_{36}\), tandis que la somme de la section prolongée utilise \(\kappa_{\mathrm{per}}\).


% BEGIN OWNER /Users/ruben/Documents/New project/output/EXTREMA_MEDIA_RAZON_RECIPROCIDAD_ES_EN_20260918/ARTICULO/ES/source/sections/k_reversibilidad.tex
\subsection{Réversibilité de la constante rationnelle et changement de phase}

La constante rationnelle \(\kappa_{\mathrm{per}}\) détermine intégralement
le registre \(K\) lorsque la base, la longueur et l'origine
de lecture sont conservées. Sa fonction numérique et sa fonction structurelle sont reliées
par une inversion exacte.

\begin{proposition}[Récupération entière et transformation cyclique]
Soit \(I(K)=\sum_{j=1}^{12}K_j1000^{12-j}\). Alors
\[
 I(K)=(1000^{12}-1)\kappa_{\mathrm{per}}.
\]
Le développement entier de \(I(K)\) en douze blocs de base \(1000\)
récupère tous les \(K_j\), y compris les zéros initiaux de chaque bloc.
Si \(\rho(K)=(K_2,\ldots,K_{12},K_1)\), on a
\begin{equation}
 \kappa(\rho K)=1000\kappa(K)-K_1.
 \label{x_reciprocidad:eq:k-rotacion-racional}
\end{equation}
\end{proposition}
\begin{proof}
La première identité inverse la définition de l'évaluation périodique.
La division euclidienne itérée par \(1000\), avec douze positions fixées,
récupère le vecteur. D'autre part,
\[
 I(\rho K)=1000I(K)-K_1(1000^{12}-1).
\]
La division par \(1000^{12}-1\) donne \eqref{x_reciprocidad:eq:k-rotacion-racional}.
\end{proof}

La récupération compose donc \(\kappa\mapsto K\mapsto U\)
et les transformations \(D_3K,D_4K,Q\) déjà construites. L'évaluation
rationnelle conserve le registre signé et les différences cycliques qui
en dépendent. La profondeur totale et les opérateurs de transport d'une
histoire restent dans leurs coordonnées propres : l'inversion précédente
a pour domaine le registre entier dodécaphasique.

Écrivons \(\kappa_j=\kappa(\rho^jK)\), avec des indices modulo \(12\).
La même identité donne
\[
 K_{j+1}=1000\kappa_j-\kappa_{j+1},\qquad
 \sum_{j=0}^{11}\kappa_j=\frac{\sum_{j=1}^{12}K_j}{999}
 =\frac{6263}{999}.
\]
La phase agit de manière exacte sur la constante : sa valeur est fixe
dans la lecture canonique, tandis que ses rotations suivent une loi affine.
La somme de l'orbite est invariante sous le choix de l'origine.

L'incidence ultérieure possède également une expression dans ces
coordonnées. Le seuil cubique détermine
\[
 B_K=\{j+1:1000\kappa_j-\kappa_{j+1}\ge729\}
 =\{4,6,9,10\}.
\]
La tétrade se récupère à partir de l'orbite rationnelle parce que celle-ci
récupère d'abord les composantes entières. Le support signé associé
à \(\alpha\) utilise en outre ses registres régionaux et la normalisation
de frontière. La relation entre coordonnée et incidence s'exprime ainsi
au moyen d'opérations identifiables.

C'est une définition opérationnelle de la fonction de \(K\) : il fournit
une mise en coordonnées entières réversible du registre orienté, dont
l'évaluation périodique et les fonctions d'incidence participent aux
constructions ultérieures. La valeur, la transformation et le domaine
sont conservés conjointement.

% END OWNER


\subsection{Périodicité et conoyau de l'opérateur de transport}
\label{x_reciprocidad:subsec:k-clase-acarreo}

La périodicité que nous venons de démontrer appartient au registre dodécaphasique. Elle n'implique
pas la périodicité des coordonnées régionales, des retenues de la clôture
ni d'alpha. Il est utile de préciser cette séparation au moyen d'un second
quotient, purement entier.

Soit \(S\) le décalage cyclique de douze coordonnées, d'orientation fixée, et soit \(b\geq2\) une base entière. L'opérateur de retenue périodique est \(\partial_b=S-bI\). Si une identité périodique s'écrit
\[
 A=P+E-\Phi-K+\partial_bC,
\]
alors les transformations
\[
 K\longmapsto K+\partial_bh,\qquad C\longmapsto C+h
\]
conservent son membre de droite. Cette invariance compare les représentants d'une
classe de retenue ; elle n'autorise pas à modifier le registre dodécaphasique déjà
sélectionné par le registre.

\begin{proposition}[Quotient de la retenue périodique]
\label{x_reciprocidad:prop:k-cociente-acarreo}
Avec l'orientation précédente,
\[
 \operatorname{coker}(S-bI)\cong\mathbb Z/(b^{12}-1)\mathbb Z.
\]
\end{proposition}

\begin{proof}
Le module cyclique se présente sous la forme
\(\mathbb Z[t]/(t^{12}-1)\). Imposer \(S=bI\) équivaut à imposer
\(t=b\) et laisse la relation \(b^{12}-1=0\). L'évaluation des
coefficients en \(b\), avec l'ordre de concaténation choisi, représente
la classe résultante.
\end{proof}

Le dénominateur \(B^{12}-1\) de \(\kappa_{\mathrm{per}}\) et ce quotient proviennent de la même longueur cyclique, mais leurs objets restent distincts : l'un est une évaluation rationnelle et l'autre un quotient d'un module de retenues. La clôture finie d'alpha sera une chaîne orientée à frontière droite. Son prolongement transportera une frontière distante compatible, et non une retenue identifiée par décret entre la première et la treizième position.

Le chapitre suivant est ainsi préparé. Le registre dodécaphasique a été fixé avant
de résoudre la récurrence d'alpha ; sa reconstruction par Hadamard est
entière et unique ; ses différences transversales restituent le registre ; et
ses deux évaluations rationnelles possèdent des domaines et des usages distincts. Aucun
de ces faits n'utilise alpha comme entrée, et aucun ne décide encore du
type arithmétique de la sortie de la clôture complète.
